\documentclass[10pt,compress,t,notes=noshow, xcolor=table]{beamer}
\newcommand{\btVFill}{\vskip0pt plus 1filll}
\newcommand\hmmax{0}
\newcommand\bmmax{0}
\input{../../style/preamble}
\input{../../latex-math/basic-math}
\input{../../latex-math/basic-ml}
\input{../../latex-math/ml-interpretable}

\usetikzlibrary{tikzmark, shapes,arrows,automata,positioning,calc,chains,trees,  shadows, decorations.pathreplacing}

\usepackage{siunitx}
% Defines macros and environments
 
\newcommand{\Xv}{\mathbf{X}} % vector X (random variable, bold)
\newcommand{\Sm}{S_m}%{Pre(\tau,j)} % S(\tau,j)
\newcommand{\Smj}{\Sm \cup \{j\}}
\newcommand{\minusSmj}{- \{\Sm \cup \{j\} \} }
\newcommand{\Stau}{S_{j}^\tau}%{Pre(\tau,j)} % S(\tau,j)
\newcommand{\Stauj}{\Stau \cup \{j\}}
\newcommand{\minusStauj}{- \{\Stau \cup \{j\} \} }


\title{Interpretable Machine Learning}
% \author{LMU}
%\institute{\href{https://compstat-lmu.github.io/lecture_iml/}{compstat-lmu.github.io/lecture\_iml}}

\begin{document}

% TODO
\titlemeta{
Shapley
}{
SHAP (SHapley Additive exPlanation)
}{
figure_man/exSHAP.png
}{
\item Recall order- and set-based definitions of Shapley values in ML
\item Interpret predictions via additive Shapley decomposition
\item Understand SHAP as surrogate-based model
\item Understand SHAP properties
}

% \newcommand{\learninggoals}{
% \item Get an intuition of additive feature attributions
% \item Understand the concept of Kernel SHAP
% \item Ability to interpret SHAP plots
% \item Global SHAP methods
% 


\begin{framev}[align=top]{Shapley Values in ML - A Short Recap}
\textbf{Shapley values (order definition):} Average over marginal contributions across all permutations of feature indices $\tau \in \Pi$:
$$
\phi_j(\xv) = \frac{1}{p!} \sum_{\tau \in \Pi} 
\underbrace{\fh_{\Stauj}(\xv_{\Stauj}) - \fh_{\Stau}(\xv_{\Stau})}_{\text{marginal contribution of feature $j$}}
$$
\begin{itemizeM}
\item For each permutation $\tau$, determine coalition $\Stau$: features before $j$ in $\tau$
\item In $\fh_S$, features not in $S$ are marginalized (e.g., randomly imputed)
%$\leadsto$ In practice, features not in \(S\) are 
\item Compute marginal contribution of adding $j$ to $\Stau$ via the difference above
\item Average over all $p!$ permutations (in practice, over $M \ll p!$)
\end{itemizeM}
\pause
\spacer[0.25]
\textbf{Alternative (set definition):} Average marginal contribution over all subsets, weighted by their relative number of appearances in permutations:
$$
\phi_j(\xv) = 
\sum_{S \subseteq \{1,\dots,p\} \setminus \{j\}} \frac{|S|!(p - |S| - 1)!}{p!} \left[\fh_{S \cup \{j\}}(\xv_{S \cup \{j\}}) - \fh_S(\xv_S)\right].
$$
\end{framev}

\begin{framev}[align=top]{Shapley Values in ML - Example}
\textbf{Example (Bike sharing data):}
\begin{itemizeM}
\item Train random forest using humidity (hum), temperature (temp), windspeed (ws)
\item Consider observation of interest $\xv$ with prediction $\fh(\xv) = \color{orange}{2573}$
\item Mean prediction $\E_{\xv}[\fh(\xv)] = \color{blue}{4515}$
\item Compute exact Shapley value for $\xv$ for feature hum:
\end{itemizeM}
\begin{center}
\begin{tabular}{c|c|c|c|c}
$S$ & $\Scupj$ & $\fh_S$ & $\fh_{\Scupj}$ & weight\\\hline
$\varnothing$& hum & \color{blue}{4515} & 4635 & $2/6$\\
temp & temp, hum & 3087 & 3060& $1/6$\\
ws & ws, hum & 4359 & 4450 & $1/6$\\
temp, ws & temp, ws, hum & 2623 & \color{orange}{2573} & $2/6$
\end{tabular}
\end{center}
\begin{itemizeM}
\item[$\Rightarrow$] {\small $\phi_{\text{hum}}(\xv) = \tfrac{2}{6}(4635-4515) + \tfrac{1}{6}(3060-3087) + \tfrac{1}{6}(4450-4359) + \tfrac{2}{6}(2573-2623) = 34$}
\item[$\Rightarrow$] {\small Analogously $\phi_{\text{temp}}(\xv) = -1654$, $\phi_{\text{ws}}(\xv) = -322$}
\end{itemizeM}
\end{framev}

\begin{framev}[align=top]{From Shapley Values to SHAP}
\splitVTT[0.51]{
%\textbf{Example (cont.):} Shapley values for other features $\phi_{\text{temp}} = -1654$, $\phi_{\text{ws}} = -322$
%\medskip
\textbf{Shapley value interpretation (for $\xv$):}
\begin{itemizeM}
\item \texttt{hum} $(+34)$ pushes pred. \emph{above} baseline (= average prediction).
\item \texttt{temp} $(-1654)$ and \texttt{ws} $(-322)$ pull prediction \emph{below} baseline.
\item Together, they explain full deviation from average prediction.
\end{itemizeM}
%\medskip
}{
\spacer[0]
%\vspace{-1cm}
\includegraphics[trim={10 17 5 5},clip, width=\linewidth]{figure/shapley2shap.pdf}
\spacer[0]
}
\pause
\textbf{Shapley-based additive decomposition} of prediction for $\xv$ gives insights on how features shift prediction from baseline $\E(\fh)$:
$$
\underbrace{{\color{orange} \fh(\xv)}}_{\text{actual prediction}}
={\color{blue} \underbrace{\phi_0}_{\E_{\mathbf X}[\fh(\mathbf X)]}}
+\sum_{j\in\{\text{hum},\text{temp},\text{ws}\}}\phi_j(\xv)
$$
%\footnotesize
$$
\phantom{=} \; {\color{orange} 2573}
= {\color{blue} 4515}
\;+\;
\bigl(34 -1654 - 322\bigr)
\;=\;4515-1942
$$
$\leadsto$ Like a LM evaluated at $\xv$: global intercept $\phi_0$ plus per-feature contribs $\phi_j(\xv)$.\\
\spacer[0.25]
%\vspace{0.4em}
\textbf{SHAP Motivation:}
Can we efficiently estimate this Shapley-based additive decomp. of $\fh(\xv)$ via a surrogate model (while preserving Shapley axioms)?
\end{framev}



\begin{framev}[align=top]{SHAP Framework \furtherreading{LUNDBERG2017}}
\textbf{SHAP} expresses the prediction of $\xv$ as a sum of contribs from each feature:
$$
g(\zv') = \phi_0 + \textstyle\sum_{j=1}^p \phi_j z'_j
$$
\begin{itemizeM}
\item $\zv' \in \{0,1\}^p$: simplified binary input referring to a coalition (coal. vector)
\item $z'_j = 1$: feature $j$ is "present" $\Rightarrow$ use $x_j$ in model evaluation
\item $z'_j = 0$: feature $j$ is "absent"\\
$\Rightarrow$ influence of $x_j$ is removed via marginalization over a reference distrib.% (averaging predictions over its possible values, given the present features)%\\
%$\leadsto$ \textit{Note: In practice (e.g., KernelSHAP), marginalization is approximated via random imputations from background data.}
\end{itemizeM}
\pause
\textbf{SHAP as a theoretical framework:} Fit a surrogate model $g(\zv')$ satisfying Shapley axioms and recovering $\fh(\xv)$ when all features are "present":
%Construct a surrogate model $g(\zv')$ such that it satisfies the Shapley value axioms and recovers $\fh(\xv)$ when all features are "present":
$$
\fh(\xv) = g(\mathbf{1}) = \phi_0 + \textstyle\sum_{j=1}^p \phi_j
$$
\textbf{Evaluation of $g(\zv')$:} Let $S = \{ j : z'_j = 1 \}$ be the active coalition. Then:
\begin{itemizeM}
\item $g(\zv') \approx \mathbb{E}[ \fh(\mathbf{X}) \mid \mathbf{X}_S = \mathbf{x}_S]\;\text{(conditional expectation)}$
\item $g(\zv') \approx \mathbb{E}_{\mathbf{X}_{-S}} [ \fh(\mathbf{x}_S, \mathbf{X}_{-S}) ] \;\;\text{(marginal expectation, i.e., PD function)}$
\item \textit{Note: Practical implementations (e.g., KernelSHAP) use the marginal expectation, approximated via random imputations from background data.}
\end{itemizeM}
\end{framev}


\begin{framev}[align=top]{SHAP Framework \furtherreading{LUNDBERG2017}}
%\textbf{Goal:} Explain a prediction \(\fh(\xv)\) by decomposing it into additive feature contributions using a local surrogate model derived from Shapley values.
\textbf{SHAP}
defines an additive surrogate $g(\zv')$ over a binary input $\zv' \in \{0,1\}^p$:
\only<1->{
\begin{exampleblock}{}
$$
g\left(\tikzmark{zp} \zv'^{(k)}\right) =
\tikzmark{ph0} \phi_0 + \sum_{j=1}^p \tikzmark{phj} \phi_j z_j'^{(k)}
$$

\begin{tikzpicture}[
remember picture,
overlay,
expl/.style={draw=blue,fill=white,rounded corners,text width=3.5cm},
arrow/.style={blue,ultra thick,->,>=latex}
]
\node<1>[expl] (zex) at (2,1.5cm)
{$\zv'^{(k)}$: \textbf{coalition vector} \\ subset of features};
\node<1>[expl] (ph0ex) at (4,-.5cm)
{$\phi_0$: \textbf{baseline} $\E[\fh(\Xv)]$};
\node<1>[expl] (phjex) at (10,0cm)
{$\phi_j$: \textbf{feature attribution} \\ marginal effect of $j$ in coalition};
\draw<1>[arrow] (zex.east) to[out=0,in=120] ([xshift= 1ex, yshift=2.3ex]{pic cs:zp});  
\draw<1>[arrow] (ph0ex.east) to[out=0,in=270] ([xshift= 1ex, yshift=-0.8ex]{pic cs:ph0});  
\draw<1>[arrow] (phjex.west) to[out=180,in=270] ([xshift= 1ex, yshift=-1ex]{pic cs:phj});
\node<2>[expl] (zex2) at (2,1cm)
{$g(\zv'^{(k)})$: \textbf{approx. prediction} for coalition};
\draw<2>[arrow] (zex2.east) to[out=0,in=140] ([xshift= 1ex, yshift=2.3ex]{pic cs:zp});  
\node<2>[expl] (phjsh) at (10,1.5cm)
{$\phi_j$: \textbf{Shapley value}};
\draw<2>[arrow] (phjsh.west) to[out=180,in=50] ([xshift= 1ex, yshift=2ex]{pic cs:phj}); 
\draw<2>[thick, decoration={brace, mirror, raise=0.5cm}, decorate]
(5.7,0.7) -- (8.2,0.7)
node[pos=0.5,below=15pt,black]{\textbf{Additive Feature Attribution}};
\end{tikzpicture}
\end{exampleblock}

\begin{onlyenv}<2>
\spacer[1]
\textbf{Next:} How do we estimate the Shapley values $\phi_j$ efficiently?
\end{onlyenv}
}
\end{framev}


%\begin{frame}{Example}

%Given the following example from the bike sharing data set

%\begin{table}[h]
%\centering
%\begin{tabular}{l rrrrr || r}
%  \hline
%  && temperature & humidity & windspeed & year & prediction\\ 
%  \hline
% example $x_{ex}$ && 24.27 & 58.5 & 13.96 & 2011 & 6825 \\ 
% \hline
%\end{tabular}
%\end{table}

%we are searching for Shapley values such that
%\begin{equation}
%\begin{array}{lllllcr}
%\phi_0 &+ \phi_{temp} &+ \phi_{hum} &+ \phi_{windspeed} &+ \phi_{yr} & = &\hat{y} \\
%4469 &+ 1809 &+ 450 &+ 241 &- 144 & = & 6825
%\end{array}
%\end{equation}

%\begin{figure}
 %   \centering
 %   \includegraphics[width=\columnwidth]{figure_man/exSHAP.png}
%\end{figure}
%\end{frame}


\begin{framev}[align=top]{Properties}
\textbf{Local Accuracy}
$$
\fh(\xv)=g\left(\mathbf{z}^{\prime}\right)=\phi_{0}+\sum_{j=1}^{p} \phi_{j} z_{j}^{\prime}
$$
\begin{onlyenv}<1>
\textbf{Intuition:} If coalition includes all features ($\mathbf{z}^{\prime} = (z^{\prime}_1, \dots, z^{\prime}_p)^\top = (1, \dots, 1)^\top$), the attributions $\phi_j$ and the baseline $\phi_0$ sum up to the original model output $\fh(\xv)$.\\
\spacer[0.25]
Local accuracy corresponds to \textbf{axiom of efficiency} in Shapley game theory
\end{onlyenv}
\begin{onlyenv}<2->
\textbf{Missingness}
$$
z_{j}^{\prime}=0 \Longrightarrow \phi_{j}=0
$$
\end{onlyenv}
\begin{onlyenv}<2>
\textbf{Intuition:} A "missing" feature (whose value is imputed) gets zero attribution
\end{onlyenv}
\begin{onlyenv}<3->
\textbf{Consistency} (Let $\mathbf{z}^{\prime (k)}_{-j} \text{ refer to } z_{j}^{\prime (k)}=0$) \\
\end{onlyenv}
\begin{onlyenv}<3>
For any two models $\fh$ and $\fh^{\prime}$, if for all inputs $\mathbf{z}^{\prime (k)} \in \{0, 1\}^p$
$$
\fh_{x}^{\prime}\left(\mathbf{z}^{\prime (k)}\right)-\fh_{x}^{\prime}\left(\mathbf{z}^{\prime (k)}_{-j}\right) \geq \fh_{x}\left(\mathbf{z}^{\prime (k)}\right)-\fh_{x}\left(\mathbf{z}^{\prime (k)} _{-j}\right) \Longrightarrow \phi_{j}\left(\fh^{\prime}, \xv\right) \geq \phi_{j}(\fh, \xv)
$$
\textbf{Intuition:} If a model changes so that the marginal contribution of a feature value increases or stays the same, the Shapley value also increases or stays the same.\\
\spacer[0.25]
\textbf{Consistency} implies Shapley’s axioms of \textbf{additivity, dummy, symmetry}.
\end{onlyenv}
\end{framev}

\endlecture
\end{document}
